\chapter{Unit 3: Analytical and Critical Thinking}


\section{Section A: Multiple Choice Practice Questions (60 Curated Problems)}

\noindent Practice and test your mastery across all core concepts of this unit. Each question reflects standard university examination patterns and Bloom's Taxonomy competencies.
\vspace{3mm}

\noindent\textbf{Q1. What is the main purpose of evaluating alternatives?} \hfill [Tag: Application of analytical reasoning and critical thinking for evaluating alternatives | Easy]
\begin{enumerate}[label=(\alph*),topsep=2pt,itemsep=1pt,partopsep=0pt,parsep=0pt]
    \item {To select the first idea}
    \item {To avoid collecting information}
    \item {To make every option identical}
    \item {To compare possible choices}
\end{enumerate}
\vspace{2mm}

\noindent\textbf{Q2. Which question best supports analytical reasoning when comparing solutions?} \hfill [Tag: Application of analytical reasoning and critical thinking for evaluating alternatives | Easy]
\begin{enumerate}[label=(\alph*),topsep=2pt,itemsep=1pt,partopsep=0pt,parsep=0pt]
    \item {Which solution has the strongest evidence?}
    \item {Which solution sounds most exciting?}
    \item {Which solution is most familiar?}
    \item {Which solution was mentioned first?}
\end{enumerate}
\vspace{2mm}

\noindent\textbf{Q3. What should be identified before comparing different alternatives?} \hfill [Tag: Application of analytical reasoning and critical thinking for evaluating alternatives | Easy]
\begin{enumerate}[label=(\alph*),topsep=2pt,itemsep=1pt,partopsep=0pt,parsep=0pt]
    \item {Random predictions}
    \item {Evaluation criteria}
    \item {Personal preferences}
    \item {Unrelated opinions}
\end{enumerate}
\vspace{2mm}

\noindent\textbf{Q4. Which action demonstrates critical thinking?} \hfill [Tag: Application of analytical reasoning and critical thinking for evaluating alternatives | Easy]
\begin{enumerate}[label=(\alph*),topsep=2pt,itemsep=1pt,partopsep=0pt,parsep=0pt]
    \item {Accepting every claim}
    \item {Checking assumptions}
    \item {Ignoring opposing views}
    \item {Choosing without reasons}
\end{enumerate}
\vspace{2mm}

\noindent\textbf{Q5. What is a cognitive bias?} \hfill [Tag: Understanding cognitive biases | Easy]
\begin{enumerate}[label=(\alph*),topsep=2pt,itemsep=1pt,partopsep=0pt,parsep=0pt]
    \item {A formal research method}
    \item {A type of numerical calculation}
    \item {A detailed project schedule}
    \item {A systematic thinking error}
\end{enumerate}
\vspace{2mm}

\noindent\textbf{Q6. Confirmation bias occurs when a person mainly seeks information that does what?} \hfill [Tag: Understanding cognitive biases | Easy]
\begin{enumerate}[label=(\alph*),topsep=2pt,itemsep=1pt,partopsep=0pt,parsep=0pt]
    \item {Challenges every belief}
    \item {Avoids all available evidence}
    \item {Supports existing beliefs}
    \item {Compares unrelated subjects}
\end{enumerate}
\vspace{2mm}

\noindent\textbf{Q7. What is the anchoring bias?} \hfill [Tag: Understanding cognitive biases | Easy]
\begin{enumerate}[label=(\alph*),topsep=2pt,itemsep=1pt,partopsep=0pt,parsep=0pt]
    \item {Using several sources equally}
    \item {Changing opinions after new evidence}
    \item {Ignoring the first available idea}
    \item {Relying too heavily on initial information}
\end{enumerate}
\vspace{2mm}

\noindent\textbf{Q8. The availability bias causes people to judge situations based mainly on what?} \hfill [Tag: Understanding cognitive biases | Easy]
\begin{enumerate}[label=(\alph*),topsep=2pt,itemsep=1pt,partopsep=0pt,parsep=0pt]
    \item {Carefully sampled data}
    \item {Official evaluation criteria}
    \item {Long-term statistical trends}
    \item {Easily recalled examples}
\end{enumerate}
\vspace{2mm}

\noindent\textbf{Q9. Which practice can help reduce the effect of cognitive biases?} \hfill [Tag: Understanding cognitive biases | Easy]
\begin{enumerate}[label=(\alph*),topsep=2pt,itemsep=1pt,partopsep=0pt,parsep=0pt]
    \item {Making decisions immediately}
    \item {Ignoring contradictory evidence}
    \item {Relying only on intuition}
    \item {Considering alternative viewpoints}
\end{enumerate}
\vspace{2mm}

\noindent\textbf{Q10. What is a rational decision-making approach?} \hfill [Tag: Decision-making approaches | Easy]
\begin{enumerate}[label=(\alph*),topsep=2pt,itemsep=1pt,partopsep=0pt,parsep=0pt]
    \item {Using goals and evidence}
    \item {Following emotions alone}
    \item {Choosing without comparison}
    \item {Copying another person's choice}
\end{enumerate}
\vspace{2mm}

\noindent\textbf{Q11. What does a satisficing approach involve?} \hfill [Tag: Decision-making approaches | Easy]
\begin{enumerate}[label=(\alph*),topsep=2pt,itemsep=1pt,partopsep=0pt,parsep=0pt]
    \item {Choosing an acceptable option}
    \item {Rejecting every available option}
    \item {Searching forever for perfection}
    \item {Selecting the most expensive option}
\end{enumerate}
\vspace{2mm}

\noindent\textbf{Q12. A decision matrix is mainly used to do what?} \hfill [Tag: Decision-making approaches | Easy]
\begin{enumerate}[label=(\alph*),topsep=2pt,itemsep=1pt,partopsep=0pt,parsep=0pt]
    \item {Predict every future event}
    \item {Compare options using criteria}
    \item {Replace all available evidence}
    \item {Record personal memories}
\end{enumerate}
\vspace{2mm}

\noindent\textbf{Q13. What is an important first step in decision-making?} \hfill [Tag: Decision-making approaches | Easy]
\begin{enumerate}[label=(\alph*),topsep=2pt,itemsep=1pt,partopsep=0pt,parsep=0pt]
    \item {Define the problem}
    \item {Avoid stating the desired outcome}
    \item {Evaluate results before acting}
    \item {Choose a solution randomly}
\end{enumerate}
\vspace{2mm}

\noindent\textbf{Q14. Which approach gives a group a structured way to reach a shared decision?} \hfill [Tag: Decision-making approaches | Easy]
\begin{enumerate}[label=(\alph*),topsep=2pt,itemsep=1pt,partopsep=0pt,parsep=0pt]
    \item {Collaborative decision-making}
    \item {Individual guessing}
    \item {Unplanned decision-making}
    \item {Automatic decision-making}
\end{enumerate}
\vspace{2mm}

\noindent\textbf{Q15. What is the purpose of a pros-and-cons list?} \hfill [Tag: Structured analysis techniques | Easy]
\begin{enumerate}[label=(\alph*),topsep=2pt,itemsep=1pt,partopsep=0pt,parsep=0pt]
    \item {To prove one opinion is correct}
    \item {To replace the decision criteria}
    \item {To organize advantages and disadvantages}
    \item {To remove all possible risks}
\end{enumerate}
\vspace{2mm}

\noindent\textbf{Q16. A SWOT analysis examines strengths, weaknesses, opportunities, and what?} \hfill [Tag: Structured analysis techniques | Easy]
\begin{enumerate}[label=(\alph*),topsep=2pt,itemsep=1pt,partopsep=0pt,parsep=0pt]
    \item {Theories}
    \item {Targets}
    \item {Timelines}
    \item {Threats}
\end{enumerate}
\vspace{2mm}

\noindent\textbf{Q17. What does a flowchart use to represent?} \hfill [Tag: Structured analysis techniques | Easy]
\begin{enumerate}[label=(\alph*),topsep=2pt,itemsep=1pt,partopsep=0pt,parsep=0pt]
    \item {Personal preferences}
    \item {Random data points}
    \item {Steps in a process}
    \item {Only financial costs}
\end{enumerate}
\vspace{2mm}

\noindent\textbf{Q18. What is the main purpose of the five whys technique?} \hfill [Tag: Structured analysis techniques | Easy]
\begin{enumerate}[label=(\alph*),topsep=2pt,itemsep=1pt,partopsep=0pt,parsep=0pt]
    \item {To estimate the final budget}
    \item {To find a root cause}
    \item {To rank team members}
    \item {To create a project logo}
\end{enumerate}
\vspace{2mm}

\noindent\textbf{Q19. What is evidence-based evaluation?} \hfill [Tag: Evidence-based evaluation methods | Easy]
\begin{enumerate}[label=(\alph*),topsep=2pt,itemsep=1pt,partopsep=0pt,parsep=0pt]
    \item {Choosing based only on popularity}
    \item {Accepting claims without questions}
    \item {Relying on untested assumptions}
    \item {Judging using reliable information}
\end{enumerate}
\vspace{2mm}

\noindent\textbf{Q20. Why is a baseline useful when evaluating a change?} \hfill [Tag: Evidence-based evaluation methods | Easy]
\begin{enumerate}[label=(\alph*),topsep=2pt,itemsep=1pt,partopsep=0pt,parsep=0pt]
    \item {It removes the need for data}
    \item {It guarantees positive results}
    \item {It identifies the most popular opinion}
    \item {It provides a comparison point}
\end{enumerate}
\vspace{2mm}

\noindent\textbf{Q21. A team is choosing between two software systems. System A has lower purchase cost, while System B has higher reliability and lower maintenance expenses. Which approach best supports a critical evaluation?} \hfill [Tag: Application of analytical reasoning and critical thinking for evaluating alternatives | Medium]
\begin{enumerate}[label=(\alph*),topsep=2pt,itemsep=1pt,partopsep=0pt,parsep=0pt]
    \item {Compare total costs and benefits over the systems' expected lifetimes}
    \item {Select the system preferred by most team members}
    \item {Choose System A because its initial cost is lower}
    \item {Choose System B because reliability is always more important}
\end{enumerate}
\vspace{2mm}

\noindent\textbf{Q22. A city is comparing three locations for a recycling center. One location has the lowest construction cost but would create longer transportation routes. Which method would provide the most balanced evaluation?} \hfill [Tag: Application of analytical reasoning and critical thinking for evaluating alternatives | Medium]
\begin{enumerate}[label=(\alph*),topsep=2pt,itemsep=1pt,partopsep=0pt,parsep=0pt]
    \item {Use weighted criteria including cost, access, and environmental impact}
    \item {Select the location supported by the largest neighborhood}
    \item {Choose the location closest to the city center}
    \item {Rank locations only by construction cost}
\end{enumerate}
\vspace{2mm}

\noindent\textbf{Q23. A proposed product feature increases customer satisfaction in a small pilot but would significantly increase production complexity. What should the team do first?} \hfill [Tag: Application of analytical reasoning and critical thinking for evaluating alternatives | Medium]
\begin{enumerate}[label=(\alph*),topsep=2pt,itemsep=1pt,partopsep=0pt,parsep=0pt]
    \item {Reject the feature because complexity is undesirable}
    \item {Ask senior managers to select the preferred option}
    \item {Launch the feature immediately}
    \item {Examine the pilot evidence and estimate broader costs and benefits}
\end{enumerate}
\vspace{2mm}

\noindent\textbf{Q24. A manager argues that a proposed training program is effective because employees who completed it received higher performance ratings. Which question most directly tests the strength of this conclusion?} \hfill [Tag: Application of analytical reasoning and critical thinking for evaluating alternatives | Medium]
\begin{enumerate}[label=(\alph*),topsep=2pt,itemsep=1pt,partopsep=0pt,parsep=0pt]
    \item {Was the training room comfortable?}
    \item {Was the program delivered during working hours?}
    \item {Did trained employees differ from others before the program began?}
    \item {Did employees enjoy the training sessions?}
\end{enumerate}
\vspace{2mm}

\noindent\textbf{Q25. After investing heavily in a failing project, a team continues funding it because members do not want the earlier investment to be wasted. Which bias is most evident?} \hfill [Tag: Understanding cognitive biases | Medium]
\begin{enumerate}[label=(\alph*),topsep=2pt,itemsep=1pt,partopsep=0pt,parsep=0pt]
    \item {Sunk cost fallacy}
    \item {Anchoring bias}
    \item {Availability bias}
    \item {Recency bias}
\end{enumerate}
\vspace{2mm}

\noindent\textbf{Q26. A hiring panel gives excessive weight to the first candidate's salary request when judging all later candidates. Which bias is affecting the evaluation?} \hfill [Tag: Understanding cognitive biases | Medium]
\begin{enumerate}[label=(\alph*),topsep=2pt,itemsep=1pt,partopsep=0pt,parsep=0pt]
    \item {Confirmation bias}
    \item {Anchoring bias}
    \item {Framing effect}
    \item {Status quo bias}
\end{enumerate}
\vspace{2mm}

\noindent\textbf{Q27. A project leader searches only for data supporting the decision already made and dismisses contradictory findings. Which practice would best reduce this bias?} \hfill [Tag: Understanding cognitive biases | Medium]
\begin{enumerate}[label=(\alph*),topsep=2pt,itemsep=1pt,partopsep=0pt,parsep=0pt]
    \item {Limit the review to the original project team}
    \item {Collect only data that supports the preferred option}
    \item {Ask a colleague to present disconfirming evidence}
    \item {Delay the decision until everyone agrees}
\end{enumerate}
\vspace{2mm}

\noindent\textbf{Q28. A product team estimates that a new service will be highly successful because a similar service recently became popular. Which bias may be influencing this prediction?} \hfill [Tag: Understanding cognitive biases | Medium]
\begin{enumerate}[label=(\alph*),topsep=2pt,itemsep=1pt,partopsep=0pt,parsep=0pt]
    \item {Sunk cost fallacy}
    \item {Availability bias}
    \item {Blind-spot bias}
    \item {Outcome bias}
\end{enumerate}
\vspace{2mm}

\noindent\textbf{Q29. Customers are more likely to choose a treatment described as having a 90\% survival rate than one described as having a 10\% mortality rate, even though the statistics are equivalent. Which bias is illustrated?} \hfill [Tag: Understanding cognitive biases | Medium]
\begin{enumerate}[label=(\alph*),topsep=2pt,itemsep=1pt,partopsep=0pt,parsep=0pt]
    \item {Authority bias}
    \item {Representativeness bias}
    \item {Planning fallacy}
    \item {Framing effect}
\end{enumerate}
\vspace{2mm}

\noindent\textbf{Q30. A team must decide whether to adopt a new technology despite uncertain demand. Which approach is most appropriate when several possible future conditions could occur?} \hfill [Tag: Decision-making approaches | Medium]
\begin{enumerate}[label=(\alph*),topsep=2pt,itemsep=1pt,partopsep=0pt,parsep=0pt]
    \item {Use scenario analysis to compare outcomes under different conditions}
    \item {Choose the option with the highest immediate popularity}
    \item {Delay all action until the future becomes completely certain}
    \item {Ignore uncertainty and use the most optimistic forecast}
\end{enumerate}
\vspace{2mm}

\noindent\textbf{Q31. A hospital must allocate limited funds among several programs. The programs have different costs, expected health benefits, and levels of risk. Which decision approach is most suitable?} \hfill [Tag: Decision-making approaches | Medium]
\begin{enumerate}[label=(\alph*),topsep=2pt,itemsep=1pt,partopsep=0pt,parsep=0pt]
    \item {First-come, first-served allocation}
    \item {Selection based only on the cheapest program}
    \item {Selection based on the loudest stakeholder}
    \item {Cost-benefit analysis with explicit risk assessment}
\end{enumerate}
\vspace{2mm}

\noindent\textbf{Q32. A manager has enough information to identify a satisfactory option but not enough time to find the absolute best option. Which decision-making concept applies?} \hfill [Tag: Decision-making approaches | Medium]
\begin{enumerate}[label=(\alph*),topsep=2pt,itemsep=1pt,partopsep=0pt,parsep=0pt]
    \item {Escalation of commitment}
    \item {Group polarization}
    \item {Satisficing}
    \item {Random choice}
\end{enumerate}
\vspace{2mm}

\noindent\textbf{Q33. A team faces a high-impact decision with uncertain outcomes and wants to reduce exposure before making a full commitment. Which approach is most appropriate?} \hfill [Tag: Decision-making approaches | Medium]
\begin{enumerate}[label=(\alph*),topsep=2pt,itemsep=1pt,partopsep=0pt,parsep=0pt]
    \item {Commit all resources to the most ambitious option}
    \item {Avoid collecting data until implementation is complete}
    \item {Select the option with the most persuasive presentation}
    \item {Implement a reversible pilot with defined evaluation criteria}
\end{enumerate}
\vspace{2mm}

\noindent\textbf{Q34. A service failure repeatedly occurs because of several interacting factors. Which technique is most useful for identifying underlying causes rather than treating symptoms?} \hfill [Tag: Structured analysis techniques | Medium]
\begin{enumerate}[label=(\alph*),topsep=2pt,itemsep=1pt,partopsep=0pt,parsep=0pt]
    \item {Root cause analysis}
    \item {Randomized brainstorming}
    \item {Preference ranking}
    \item {Simple majority voting}
\end{enumerate}
\vspace{2mm}

\noindent\textbf{Q35. A team wants to compare four solutions using criteria that have different levels of importance. Which tool would best organize this comparison?} \hfill [Tag: Structured analysis techniques | Medium]
\begin{enumerate}[label=(\alph*),topsep=2pt,itemsep=1pt,partopsep=0pt,parsep=0pt]
    \item {Open-ended opinion poll}
    \item {Decision matrix with weighted criteria}
    \item {Unstructured discussion}
    \item {Single-factor ranking}
\end{enumerate}
\vspace{2mm}

\noindent\textbf{Q36. A project manager wants to identify tasks that could delay the entire project if they take longer than planned. Which technique should be used?} \hfill [Tag: Structured analysis techniques | Medium]
\begin{enumerate}[label=(\alph*),topsep=2pt,itemsep=1pt,partopsep=0pt,parsep=0pt]
    \item {Stakeholder mapping}
    \item {Affinity diagramming}
    \item {PESTLE analysis}
    \item {Critical path analysis}
\end{enumerate}
\vspace{2mm}

\noindent\textbf{Q37. A company is exploring whether to enter a new market. It wants to examine internal capabilities and external opportunities and threats. Which framework is most appropriate?} \hfill [Tag: Structured analysis techniques | Medium]
\begin{enumerate}[label=(\alph*),topsep=2pt,itemsep=1pt,partopsep=0pt,parsep=0pt]
    \item {Five whys}
    \item {Pareto chart}
    \item {SWOT analysis}
    \item {Process mapping}
\end{enumerate}
\vspace{2mm}

\noindent\textbf{Q38. A school introduces a new teaching method in one class and compares its results with a similar class using the existing method. What makes this evaluation stronger?} \hfill [Tag: Evidence-based evaluation methods | Medium]
\begin{enumerate}[label=(\alph*),topsep=2pt,itemsep=1pt,partopsep=0pt,parsep=0pt]
    \item {Evaluating the method solely through student opinions}
    \item {Using a comparison group with similar baseline characteristics}
    \item {Measuring only the class using the new method}
    \item {Allowing each teacher to define success differently}
\end{enumerate}
\vspace{2mm}

\noindent\textbf{Q39. A survey about workplace satisfaction receives responses from only 12\% of employees, mostly from people with strong opinions. Which concern is most important?} \hfill [Tag: Evidence-based evaluation methods | Medium]
\begin{enumerate}[label=(\alph*),topsep=2pt,itemsep=1pt,partopsep=0pt,parsep=0pt]
    \item {The results may suffer from response bias}
    \item {The findings cannot contain any useful information}
    \item {The sample is automatically representative}
    \item {The survey necessarily proves dissatisfaction}
\end{enumerate}
\vspace{2mm}

\noindent\textbf{Q40. A company reports that sales rose after a new advertising campaign. Which additional evidence would best support the claim that the campaign caused the increase?} \hfill [Tag: Evidence-based evaluation methods | Medium]
\begin{enumerate}[label=(\alph*),topsep=2pt,itemsep=1pt,partopsep=0pt,parsep=0pt]
    \item {The advertising agency predicted higher sales}
    \item {The campaign received positive comments from employees}
    \item {Sales also rose in regions where the campaign was not used}
    \item {A similar control region showed no comparable increase}
\end{enumerate}
\vspace{2mm}

\noindent\textbf{Q41. A city must choose between projects A and B. A has a higher expected benefit, but its benefits vary substantially across scenarios; B has a lower expected benefit but performs consistently. Which analysis best supports a defensible choice?} \hfill [Tag: Application of analytical reasoning and critical thinking for evaluating alternatives | Hard]
\begin{enumerate}[label=(\alph*),topsep=2pt,itemsep=1pt,partopsep=0pt,parsep=0pt]
    \item {Select the project with the larger best-case benefit}
    \item {Compare scenario-weighted outcomes and risk preferences}
    \item {Select B because variability always indicates poor value}
    \item {Select A because its mean benefit is higher}
\end{enumerate}
\vspace{2mm}

\noindent\textbf{Q42. An intervention appears superior to its alternatives because it improves three measured outcomes but worsens an unmeasured safety outcome. What is the strongest analytical response?} \hfill [Tag: Application of analytical reasoning and critical thinking for evaluating alternatives | Hard]
\begin{enumerate}[label=(\alph*),topsep=2pt,itemsep=1pt,partopsep=0pt,parsep=0pt]
    \item {Rank it first using the three measured outcomes}
    \item {Delay comparison until relevant outcomes are included}
    \item {Average all available outcomes without assigning weights}
    \item {Exclude safety because it lacks numerical evidence}
\end{enumerate}
\vspace{2mm}

\noindent\textbf{Q43. A proposed solution has a lower total cost than its alternatives, but the estimate assumes unusually favorable maintenance conditions. Which test most directly examines the robustness of the recommendation?} \hfill [Tag: Application of analytical reasoning and critical thinking for evaluating alternatives | Hard]
\begin{enumerate}[label=(\alph*),topsep=2pt,itemsep=1pt,partopsep=0pt,parsep=0pt]
    \item {Conduct sensitivity analysis on maintenance assumptions}
    \item {Replace the estimate with the highest historical cost}
    \item {Choose the alternative with the lowest initial price}
    \item {Remove maintenance from every alternative's calculation}
\end{enumerate}
\vspace{2mm}

\noindent\textbf{Q44. Two alternatives have equal expected utility, but one preserves the option to scale up after learning from a pilot. Which principle favors that alternative?} \hfill [Tag: Application of analytical reasoning and critical thinking for evaluating alternatives | Hard]
\begin{enumerate}[label=(\alph*),topsep=2pt,itemsep=1pt,partopsep=0pt,parsep=0pt]
    \item {Sunk-cost minimization}
    \item {Majority preference aggregation}
    \item {Option value under uncertainty}
    \item {Status-quo preservation}
\end{enumerate}
\vspace{2mm}

\noindent\textbf{Q45. A team continues funding a failing project because it has already spent two years developing it. Which corrective practice most directly addresses the bias involved?} \hfill [Tag: Understanding cognitive biases | Hard]
\begin{enumerate}[label=(\alph*),topsep=2pt,itemsep=1pt,partopsep=0pt,parsep=0pt]
    \item {Include prior expenditure in future benefit estimates}
    \item {Evaluate continuation using future costs and benefits}
    \item {Increase funding to protect earlier investments}
    \item {Ask the original sponsor to defend the project}
\end{enumerate}
\vspace{2mm}

\noindent\textbf{Q46. After observing several successful outcomes from a risky strategy, a manager concludes that the strategy is reliable despite a small sample. Which bias is most evident?} \hfill [Tag: Understanding cognitive biases | Hard]
\begin{enumerate}[label=(\alph*),topsep=2pt,itemsep=1pt,partopsep=0pt,parsep=0pt]
    \item {Availability and representativeness bias}
    \item {Reactance against external criticism}
    \item {Loss aversion toward failed strategies}
    \item {Anchoring on the original forecast}
\end{enumerate}
\vspace{2mm}

\noindent\textbf{Q47. A negotiator receives an initial price of 1.2 million and makes all later estimates close to that figure, although independent market data suggest a much lower value. Which debiasing method is most appropriate?} \hfill [Tag: Understanding cognitive biases | Hard]
\begin{enumerate}[label=(\alph*),topsep=2pt,itemsep=1pt,partopsep=0pt,parsep=0pt]
    \item {Average the anchor with the final proposal}
    \item {Use the anchor as a conservative planning baseline}
    \item {Ask whether the initial figure feels realistic}
    \item {Generate an independent estimate before reviewing the anchor}
\end{enumerate}
\vspace{2mm}

\noindent\textbf{Q48. A project team seeks only studies supporting its preferred design and dismisses contradictory findings as irrelevant. Which intervention best reduces this pattern?} \hfill [Tag: Understanding cognitive biases | Hard]
\begin{enumerate}[label=(\alph*),topsep=2pt,itemsep=1pt,partopsep=0pt,parsep=0pt]
    \item {Delay criticism until implementation begins}
    \item {Define disconfirming evidence before reviewing results}
    \item {Invite more supporters of the preferred design}
    \item {Use the most frequently cited supporting study}
\end{enumerate}
\vspace{2mm}

\noindent\textbf{Q49. A forecast is repeatedly too optimistic, yet each new forecast is adjusted only slightly because managers believe the next case will be different. Which bias should be investigated first?} \hfill [Tag: Understanding cognitive biases | Hard]
\begin{enumerate}[label=(\alph*),topsep=2pt,itemsep=1pt,partopsep=0pt,parsep=0pt]
    \item {Hindsight bias caused by retrospective analysis}
    \item {Planning fallacy reinforced by inside-view reasoning}
    \item {Omission bias caused by excessive caution}
    \item {Framing effect caused by neutral wording}
\end{enumerate}
\vspace{2mm}

\noindent\textbf{Q50. A decision has irreversible consequences, uncertain probabilities, and potentially catastrophic losses. Which approach is most defensible?} \hfill [Tag: Decision-making approaches | Hard]
\begin{enumerate}[label=(\alph*),topsep=2pt,itemsep=1pt,partopsep=0pt,parsep=0pt]
    \item {Use equal probabilities for every possible outcome}
    \item {Choose the option with the highest average payoff}
    \item {Apply a precautionary approach with staged commitments}
    \item {Select the most popular option among stakeholders}
\end{enumerate}
\vspace{2mm}

\noindent\textbf{Q51. A group must decide whether to launch a product. Evidence is incomplete, but delaying has a measurable opportunity cost. Which decision rule best reflects this trade-off?} \hfill [Tag: Decision-making approaches | Hard]
\begin{enumerate}[label=(\alph*),topsep=2pt,itemsep=1pt,partopsep=0pt,parsep=0pt]
    \item {Wait until uncertainty is completely eliminated}
    \item {Compare the value of information with delay costs}
    \item {Launch automatically because delay has a cost}
    \item {Choose according to the executive's confidence}
\end{enumerate}
\vspace{2mm}

\noindent\textbf{Q52. A committee must choose among alternatives with conflicting stakeholder values that cannot be reduced to one common metric. Which approach is most suitable?} \hfill [Tag: Decision-making approaches | Hard]
\begin{enumerate}[label=(\alph*),topsep=2pt,itemsep=1pt,partopsep=0pt,parsep=0pt]
    \item {Use transparent multi-criteria deliberation}
    \item {Force every criterion into monetary units}
    \item {Select the option with the highest single score}
    \item {Let the largest stakeholder group decide}
\end{enumerate}
\vspace{2mm}

\noindent\textbf{Q53. A decision maker chooses the first solution that meets minimum requirements, even though better solutions may exist and search costs are substantial. Which approach is being used?} \hfill [Tag: Decision-making approaches | Hard]
\begin{enumerate}[label=(\alph*),topsep=2pt,itemsep=1pt,partopsep=0pt,parsep=0pt]
    \item {Maximization under complete information}
    \item {Satisficing under bounded rationality}
    \item {Regret minimization under certainty}
    \item {Randomized choice under ambiguity}
\end{enumerate}
\vspace{2mm}

\noindent\textbf{Q54. Two policies have similar expected outcomes, but policy X could create severe losses in a low-probability scenario while policy Y cannot. Under a minimax-regret approach, what should be examined?} \hfill [Tag: Decision-making approaches | Hard]
\begin{enumerate}[label=(\alph*),topsep=2pt,itemsep=1pt,partopsep=0pt,parsep=0pt]
    \item {The policy favored by the largest coalition}
    \item {The worst regret associated with each policy}
    \item {The average result across all scenarios}
    \item {The probability of the most favorable scenario}
\end{enumerate}
\vspace{2mm}

\noindent\textbf{Q55. A causal diagram shows that training improves performance, but only when employees have access to functioning equipment. Which structure does this relationship represent?} \hfill [Tag: Structured analysis techniques | Hard]
\begin{enumerate}[label=(\alph*),topsep=2pt,itemsep=1pt,partopsep=0pt,parsep=0pt]
    \item {A moderator that changes effect strength}
    \item {A mediator that transmits every effect}
    \item {A collider caused by both measured variables}
    \item {A confounder unrelated to the intervention}
\end{enumerate}
\vspace{2mm}

\noindent\textbf{Q56. In a fault-tree analysis, system failure occurs if either subsystem A fails or both subsystems B and C fail together. Which logical expression represents this structure?} \hfill [Tag: Structured analysis techniques | Hard]
\begin{enumerate}[label=(\alph*),topsep=2pt,itemsep=1pt,partopsep=0pt,parsep=0pt]
    \item {( A $\lor$}
    \item {A $\lor$}
    \item {( A $\land$}
    \item {A $\land$}
\end{enumerate}
\vspace{2mm}

\noindent\textbf{Q57. A team maps a problem as symptoms, immediate causes, and deeper systemic causes. Several symptoms trace back to the same policy constraint. What is the main analytical benefit of this map?} \hfill [Tag: Structured analysis techniques | Hard]
\begin{enumerate}[label=(\alph*),topsep=2pt,itemsep=1pt,partopsep=0pt,parsep=0pt]
    \item {It guarantees that one intervention solves every symptom}
    \item {It replaces the need to validate causal assumptions}
    \item {It ranks causes according to stakeholder seniority}
    \item {It distinguishes leverage points from visible effects}
\end{enumerate}
\vspace{2mm}

\noindent\textbf{Q58. A program's participants improve after implementation, but there is no comparison group and outcomes may reflect seasonal trends. Which conclusion is most justified?} \hfill [Tag: Evidence-based evaluation methods | Hard]
\begin{enumerate}[label=(\alph*),topsep=2pt,itemsep=1pt,partopsep=0pt,parsep=0pt]
    \item {The program had no effect because confounding exists}
    \item {The improvement proves the program is cost-effective}
    \item {The change is consistent with benefit but not causal proof}
    \item {The program caused the observed improvement}
\end{enumerate}
\vspace{2mm}

\noindent\textbf{Q59. A randomized trial finds a statistically significant effect that is much smaller than the threshold users consider meaningful. Which interpretation is best?} \hfill [Tag: Evidence-based evaluation methods | Hard]
\begin{enumerate}[label=(\alph*),topsep=2pt,itemsep=1pt,partopsep=0pt,parsep=0pt]
    \item {The sample must be invalid because significance was detected}
    \item {Statistical significance establishes practical importance}
    \item {The result may be real but practically unimportant}
    \item {The intervention is ineffective because the effect is small}
\end{enumerate}
\vspace{2mm}

\noindent\textbf{Q60. A review combines studies with highly different populations, interventions, and outcome measures, producing a single pooled estimate. What should analysts assess before trusting that estimate?} \hfill [Tag: Evidence-based evaluation methods | Hard]
\begin{enumerate}[label=(\alph*),topsep=2pt,itemsep=1pt,partopsep=0pt,parsep=0pt]
    \item {Whether heterogeneity undermines meaningful pooling}
    \item {Whether the oldest study has the largest weight}
    \item {Whether the pooled estimate is numerically positive}
    \item {Whether every study has an identical sample size}
\end{enumerate}
\vspace{2mm}


\newpage
\section{Section B: Comprehensive Subjective Problem with Detailed Analytical Rubric}

\begin{subjectiveq}{3}{1}
\textbf{Case Background:} An educational technology university campus is debating replacing all traditional physical student identity cards and access turnstiles with an AI-driven biometric facial recognition system. The proposal has caused division among university leadership, students, and cybersecurity personnel.

\textbf{Required Design Deliverables:}
\begin{enumerate}[label=(\alph*)]
    \item Identify and analyze three cognitive biases (e.g., Automation Bias, Confirmation Bias, Sunk Cost Fallacy) likely distorting stakeholder evaluation.
    \item Execute a Multi-Criteria Decision Analysis (MCDA) evaluating three alternative solutions against four weighted criteria.
    \item Formulate a bounded rationality risk mitigation framework applying the MECE (Mutually Exclusive, Collectively Exhaustive) principle.
\end{enumerate}
\end{subjectiveq}

\begin{solbox}
\textbf{1. Cognitive Bias Audit \& Analytical De-biasing:}\\
\begin{itemize}
    \item \textbf{Automation Bias / Novelty Bias:} University IT administrators assume that because facial recognition involves cutting-edge AI, it is inherently superior to simple RFID tokens, overlooking lighting variations and false match rates.
    \item \textbf{Confirmation Bias:} Security chiefs selectively seek vendor whitepapers demonstrating 99.8\% lab accuracy while dismissing independent audits demonstrating algorithmic demographic disparity.
    \item \textbf{Sunk Cost Fallacy:} Administrators emphasize that \$80,000 was already spent on biometric consulting, justifying continued investment despite mounting student privacy protests.
\end{itemize}

\textbf{2. Multi-Criteria Decision Analysis (MCDA) Evaluation:}\\
Criteria Weights: $W_1$ Security \& Data Privacy (0.35); $W_2$ Speed \& Throughput (0.25); $W_3$ Capital \& Operating Cost (0.20); $W_4$ User Consent \& Inclusion (0.20). [Scale 1 to 10]

{\small
\begin{tabularx}{\linewidth}{|X|c|c|c|c|c|} \hline
\textbf{Alternative Candidate} & \textbf{Privacy} & \textbf{Speed} & \textbf{Cost} & \textbf{Consent} & \textbf{Weighted} \\
 & \textbf{(0.35)} & \textbf{(0.25)} & \textbf{(0.20)} & \textbf{(0.20)} & \textbf{Score} \\ \hline
Option A: Full Facial Recognition & 4.0 (1.40) & 9.0 (2.25) & 3.0 (0.60) & 4.0 (0.80) & \textbf{5.05} \\ \hline
Option B: NFC Mobile Wallet Card & 8.5 (2.97) & 8.5 (2.12) & 7.5 (1.50) & 9.0 (1.80) & \textbf{8.40} \\ \hline
Option C: Legacy Plastic Barcode & 6.0 (2.10) & 5.0 (1.25) & 9.0 (1.80) & 8.0 (1.60) & \textbf{6.75} \\ \hline
\end{tabularx}\par}

\textit{Decision Conclusion:} Option B (Encrypted NFC Mobile Credentials) emerges as the clear winner with a weighted composite score of 8.40, outperforming AI facial recognition by 66\% due to vastly superior privacy preservation and user consent.

\textbf{3. MECE Risk Assessment Architecture:}\\
To ensure exhaustive coverage without conceptual overlap, risks are divided across four MECE pillars:
\begin{enumerate}[label=(\roman*)]
    \item \textbf{Technical \& Infrastructure:} Battery depletion on student phones; offline gateway validation failure; zero-day BLE/NFC protocol vulnerabilities.
    \item \textbf{Operational \& Logistical:} Onboarding delays during freshman orientation week; replacement latency for lost phones.
    \item \textbf{Legal \& Regulatory:} Compliance with national digital personal data protection frameworks; biometric cross-border cloud storage restrictions.
    \item \textbf{Sociological \& Behavioral:} Accessibility accommodations for students without modern smartphones through subsidized backup smart-fobs.
\end{enumerate}
\end{solbox}

\begin{examinertip}
\textbf{Criteria for Full 10/10 Marks:}\\
1. \textbf{Rigorous Theoretical Framing (3 Marks):} Explicitly identify and define the formal model or framework applied (e.g., d.school stages, 5-Whys, MCDA, SCAMPER, Blueprinting, Minto SCQA). Avoid vague hand-waving.\\
2. \textbf{Contextual Application (4 Marks):} Directly connect every step to the specific facts, users, and constraints given in the case scenario.\\
3. \textbf{Actionable Deliverables \& Invariants (3 Marks):} Provide precise, formatted tables, matrices, or SMART criteria with clear quantitative bounds.
\end{examinertip}
